\subsection{向量场}

设 X 是类为 \(C^{r}\) 的流形，且 U 是 X 的一个开集。

8.2.1. 切丛 T(X) 在 U 上的一个截面（不必连续，参见记号和约定）称为 U 上的一个向量场。

8.2.2. T'(X)（7.7.3）是 T(X) 的对偶丛，称为 X 的余切丛。它的截面称为余向量场。若 f 是 X 上取值于 K 的类为 \(C^{k}\) 的函数（其中 \(1 \leqslant k \leqslant r\)），则其微分 df: \(x \mapsto d_{x}f\)（5.5.6）是类为 \(C^{k-1}\) 的余向量场。更一般地，若 f 是 X 上取值于巴拿赫空间 E 的类为 \(C^{k}\) 的函数（ \(1 \leqslant k \leqslant r\)），则其微分 df 是类为 \(C^{k-1}\) 的丛 \(\mathcal{L}(\mathrm{T}(\mathrm{X}); \mathrm{E}_{\mathrm{X}})\) 的截面。

8.2.3. 设 \(\xi\) 是 X 上类为 \(\mathbf{C}^{r-1}\) 的向量场，且 \(f\in\mathcal{C}^{r}(X;F)\) 是一个取值于巴拿赫空间 F 的类为 \(\mathbf{C}^{r}\) 的函数。将函数 \(\langle\xi,df\rangle\)、或 \(\mathbf{D}_{\xi}(f)\)、或仍记为 \(\xi(f)\) 记作 \(x\mapsto d_{x}f(\xi(x))\)；它属于 \(\mathcal{C}^{r-1}(X;F)\)。对于固定的 \(f\)，映射 \(\xi\mapsto\mathbf{D}_{\xi}(f)\) 是 \(\mathcal{C}^{r-1}(X)\)-线性的。若 \(\mathbf{D}_{\xi}f=0\)，对于定义在 X 的某个开集上的、取值于某个巴拿赫空间的每个函数 \(f\)，其类为 \(\mathbf{C}^{r}\)，则 \(\xi=0\)。

设 E、F、G 是巴拿赫空间，且 \((u, v) \mapsto u \cdot v\) 是从 \(\mathbf{E} \times \mathbf{F}\) 到 \(\mathbf{G}\) 的连续双线性映射。设 \(f \in \mathcal{C}^r(\mathbf{X}; \mathbf{E})\)、\(g \in \mathcal{C}^r(\mathbf{X}; \mathbf{F})\)，并令 \(f.g \in \mathcal{C}^r(\mathbf{X}; \mathbf{G})\) 为它们的乘积。若 \(\xi\) 是类为 \(\mathbf{C}^{r-1}\) 的 \(\mathbf{X}\) 上的向量场，则有

\[
\mathrm{D} _ {\xi} (f. g) = \mathrm{D} _ {\xi} (f). g + f. \mathrm{D} _ {\xi} (g).
\]

特别地，映射 \(\mathbf{D}_{\xi}\colon\mathcal{C}^{r}(\mathbf{X})\to\mathcal{C}^{r-1}(\mathbf{X})\) 是从代数 \(\mathcal{C}^{r}(\mathbf{X})\) 到 \(\mathcal{C}^{r}(\mathbf{X})\)-模 \(\mathcal{C}^{r-1}(\mathbf{X})\) 的一个导子。

8.2.4. 假设 X 是局部有限维的，并且下列两个假设之一成立：

\begin{enumerate}
\def\labelenumi{(\roman{enumi})}
\item
  \(r = \infty\) 且 X 是豪斯多夫的；
\item
  \(r = \omega\) 且 X 同构于 \(K^{n}\) 的一个开多圆盘。
\end{enumerate}

于是，对于代数 \(\mathcal{C}^r(X)\) 的每个导子 D，存在唯一一个向量场 \(\xi\)，其类为 \(C^r\)，使得 \(D = D_{\xi}\)。

8.2.5. 设 \((u^{1},\ldots ,u^{n})\) 是 U 上 X 的一个坐标系，且 \(\xi\) 是 U 上的一个向量场。第 \(i\) 个坐标是 \(\xi\) 在由 \((u^{1},\ldots ,u^{n})\) 定义的切标架（8.1.5）中的坐标，等于 \(\mathrm{D}_{\xi}(u^i)\)。换言之，有

\[
\xi = \sum_ {1 \leqslant i \leqslant n} \mathrm{D} _ {\xi} (u ^ {i}). \partial / \partial u ^ {i}.
\]

8.2.6. 设 \(\varphi: X \to Y\) 是类为 \(C^r\) 的流形态射，\(\xi\) 是 \(X\) 上的一个向量场，\(\eta\) 是 \(Y\) 上的一个向量场。称 \(\xi\) 关于 \(\varphi\) 与 \(\eta\) 相关，如果对所有 \(x \in X\) 都有 \(\eta(\varphi(x)) = T_x(\varphi)(\xi(x))\)。若情况如此，且 \(\xi\) 和 \(\eta\) 都是类为 \(C^{r-1}\) 的，则对于每个巴拿赫空间 E，图表

\[
\begin{array}{c c c} \mathcal {C} ^ {r} (\mathrm{Y}; \mathrm{E}) & \xrightarrow {\varphi^ {*}} & \mathcal {C} ^ {r} (\mathrm{X}; \mathrm{E}) \\ \Big \downarrow_ {\mathrm{D} _ {\eta}} & & \Big \downarrow_ {\mathrm{D} _ {\xi}} \\ \mathcal {C} ^ {r - 1} (\mathrm{Y}; \mathrm{E}) & \xrightarrow {\varphi^ {*}} & \mathcal {C} ^ {r - 1} (\mathrm{X}; \mathrm{E}) \end{array}
\]

\(\mathrm{where}\ \varphi^{*}(f) = f\circ \varphi ,\) 是交换的。

当 \(\varphi\) 是平展时，对于 Y 上的每个向量场 \(\eta\)，存在唯一一个 X 上的向量场，它关于 \(\varphi\) 与 \(\eta\) 相关；记为 \(\varphi^{*}\eta\)。

8.2.7. 设 \(\varphi: X \to Y\) 是类为 \(C^r\) 的流形态射。如果 \(\omega\) 是一个余向量场 \(^{(1)}\)，定义在 \(Y\) 上，则定义 \(\varphi^*\omega\)，它是 \(X\) 上的余向量场，如下：

\[
(\varphi^ {*} \omega) (x) = ^ {t} (\mathrm{T} _ {x} (\varphi)) (\omega (\varphi (x))) (x \in X).
\]

更一般地，设 \(\tau\) 是一个类为 \(\mathbf{C}^{r-1}\) 的向量函子（7.6），其类型为 \((\mathbf{I}_{+}, \mathbf{I}_{-})\)，其中 \(\mathbf{I}_{+} = \varnothing\)，且 \(\mathbf{I}_{-}\) 约化为单个元素（这样的函子称为逆变的）。若 \(\omega\) 是类为 \(\mathbf{C}^{r-1}\) 的向量丛 \(\tau(\mathrm{T}(\mathrm{Y}))\) 的、以 \(\mathbf{Y}\) 为基底的截面，则定义 \(\varphi^{*}\omega\) 为丛 \(\tau(\mathrm{T}(\mathrm{X}))\) 的截面：\((\varphi^{*}\omega)(x) = \tau(\mathrm{T}_{x}(\varphi))(\omega(\varphi(x)))\)（ \(x \in \mathrm{X}\)）。
% semantic-fragments: legacy-input-seam
\relax{}
